% Pollis Data Packages: the datasets of the examples (PollisDatasets.jl) -- draft 1 (28 October 2026)
% Compile with: tectonic -X compile PollisDatasets.tex
% Shared look for the Pollis New Methods documents (same fonts and colours as the brochure and the workflow handouts).
\documentclass[10pt,a4paper]{article}

\usepackage[a4paper,left=18mm,right=18mm,top=15mm,bottom=17mm,footskip=8mm]{geometry}
\usepackage[T1]{fontenc}
\usepackage[default,semibold]{sourcesanspro}
\usepackage{textcomp}
\usepackage[italic]{mathastext}
\usepackage{amsmath}
\usepackage{microtype}
\usepackage[table]{xcolor}
\usepackage{tikz}
\usepackage[most]{tcolorbox}
\usepackage{fontawesome5}
\usepackage{tabularx,booktabs}
\usepackage{enumitem}
\usepackage{fancyhdr}
\usepackage[hidelinks]{hyperref}

% ---------------------------------------------------------------- colours (as in the brochure)
\definecolor{navy}{HTML}{1C3A5E}
\definecolor{pblue}{HTML}{3B78B5}
\definecolor{amber}{HTML}{F2B33D}
\definecolor{deepamber}{HTML}{D9902A}
\definecolor{pale}{HTML}{EEF3F9}
\definecolor{ink}{HTML}{2A2F36}
\definecolor{logonavy}{HTML}{173D63}
\definecolor{logodark}{HTML}{C98A1E}
\definecolor{logoamber}{HTML}{F2B446}
\definecolor{logocream}{HTML}{FEEFD2}
\definecolor{soft}{HTML}{5E6773}
\definecolor{agree}{HTML}{2F7D4F}
\definecolor{differ}{HTML}{B23B3B}
\color{ink}

\setlength{\parindent}{0pt}
\setlength{\parskip}{4pt}
\setlist{nosep,leftmargin=*,itemsep=2pt}
\renewcommand{\arraystretch}{1.15}

\pagestyle{fancy}
\fancyhf{}
\renewcommand{\headrulewidth}{0pt}
\fancyfoot[L]{\footnotesize\color{soft}\textbf{Pollis}\enspace New Methods\enspace·\enspace\nmtheme}
\fancyfoot[R]{\footnotesize\color{soft}\thepage}
\newcommand{\nmtheme}{}
\newcommand{\nmdate}{11 October 2026}

% The Pollis logo, redrawn so it prints sharply at any size
\newcommand{\pollislogo}[1]{%
  \begin{tikzpicture}[x=#1/512,y=-#1/512,baseline=(current bounding box.center)]
    \fill[logonavy] (0,0) rectangle (512,512);
    \fill[logodark] (36,422) -- (104,256) -- (172,422) -- cycle;
    \fill[logoamber] (104,422) -- (203,180) -- (302,422) -- cycle;
    \fill[logocream] (203,422) -- (339,89) -- (476,422) -- cycle;
  \end{tikzpicture}}

% Top of the first page: logo, name; then the kicker, title and a one-line promise
\newcommand{\nmhead}[3]{%
  \renewcommand{\nmtheme}{#1}%
  \thispagestyle{fancy}%
  \begin{minipage}[c]{13mm}\pollislogo{11mm}\end{minipage}%
  \begin{minipage}[c]{\dimexpr\linewidth-13mm\relax}
    {\color{navy}\Large\bfseries Pollis}\enspace{\color{soft}New Methods}\hfill
    {\color{deepamber}\small\bfseries\MakeUppercase{#1}}
  \end{minipage}\par\vspace{7mm}
  {\color{navy}\fontsize{22}{26}\selectfont\bfseries #2}\par\vspace{4pt}
  {\color{soft}\large #3}\par\vspace{5mm}}

\newcommand{\nmsection}[1]{\par\vspace{5pt}{\color{navy}\large\bfseries #1}\par\vspace{1pt}}
\newcommand{\nmsub}[1]{\par\vspace{2pt}{\color{deepamber}\bfseries #1}\par}
\newcommand{\gt}{\,{\color{deepamber}\bfseries ›}\,}
\newcommand{\menupath}[1]{{\bfseries\color{navy}#1}}
\newcommand{\fn}[1]{\texttt{\small #1}}
% Citations: \nmcite{key}{Breiman (2001)} links to the entry \nmref{key}{...} in the References
\newcommand{\nmcite}[2]{\hyperlink{ref:#1}{\color{pblue}#2}}
\newcommand{\nmref}[2]{\raisebox{1.1em}[0pt][0pt]{\hypertarget{ref:#1}{}}#2\par}
\newcommand{\nmdoi}[1]{\href{https://doi.org/#1}{\color{pblue}doi:#1}}
% Where a notebook is in Pollis: its title in Browse Methods, and its file
\newcommand{\browse}[2]{\menupath{Compose \gt New Methods \gt Browse Methods}: \emph{#1} (\fn{#2.ipynb})}

% Key facts at the top: source, missing in Julia, notebook, result
\newcommand{\factsheet}[4]{%
  \begin{tcolorbox}[enhanced, colback=pale, colframe=pale, boxrule=0pt, sharp corners,
    left=7pt, right=7pt, top=5pt, bottom=5pt, fontupper=\small]
  \begin{tabularx}{\linewidth}{@{}>{\leavevmode\bfseries\color{navy}}p{27mm}>{\leavevmode}X@{}}
  Tested against & #1\\
  Missing in Julia & #2\\
  Notebook & #3\\
  Result & #4
  \end{tabularx}
  \end{tcolorbox}}

% Boxes
\newtcolorbox{pbox}[2][]{enhanced, colback=pale, colframe=amber, boxrule=0pt,
  leftrule=3pt, sharp corners, left=7pt, right=7pt, top=6pt, bottom=6pt,
  fontupper=\small\raggedright, before upper={{\color{navy}\normalsize\bfseries #2}\par\smallskip}, #1}
\newtcolorbox{nbox}[2][]{enhanced, colback=navy, colframe=navy, boxrule=0pt,
  sharp corners, left=8pt, right=8pt, top=7pt, bottom=7pt, coltext=white,
  fontupper=\small\raggedright, before upper={{\color{amber}\normalsize\bfseries #2}\par\smallskip}, #1}

% Value tables: published beside Julia
\newenvironment{valuetable}[1]{%
  \par\small\tabularx{\linewidth}{#1}\toprule}{\bottomrule\endtabularx\par}
\newcommand{\ok}{{\color{agree}\faCheck}}
\newcommand{\no}{{\color{differ}\faTimes}}

% End of the document: credit and licence
\newcommand{\nmcredit}{%
  \par\vfill
  \begin{minipage}[b]{0.78\linewidth}\footnotesize\color{soft}
  Methods written and tested with Claude Code (Anthropic) in Pollis, from the published papers and their data.\par
  \textcopyright\ 2026 Antonio Saragga Seabra\enspace·\enspace Licensed under \href{https://creativecommons.org/licenses/by-sa/4.0/}{CC BY-SA 4.0}\enspace·\enspace Draft, \nmdate
  \end{minipage}\hfill
  \begin{minipage}[b]{0.18\linewidth}\raggedleft\pollislogo{9mm}\par\vspace{2pt}{\footnotesize contact@pollis.org}\end{minipage}}


% A data package Pollis installs for its examples, not one that reproduces a paper: its own header, key facts and credit
\renewcommand{\nmdate}{28 October 2026}
\fancyfoot[L]{\footnotesize\color{soft}\textbf{Pollis}\enspace Data Packages\enspace·\enspace\nmtheme}
\renewcommand{\nmhead}[3]{%
  \renewcommand{\nmtheme}{#1}%
  \thispagestyle{fancy}%
  \begin{minipage}[c]{13mm}\pollislogo{11mm}\end{minipage}%
  \begin{minipage}[c]{\dimexpr\linewidth-13mm\relax}
    {\color{navy}\Large\bfseries Pollis}\enspace{\color{soft}Data Packages}\hfill
    {\color{deepamber}\small\bfseries\MakeUppercase{#1}}
  \end{minipage}\par\vspace{7mm}
  {\color{navy}\fontsize{22}{26}\selectfont\bfseries #2}\par\vspace{4pt}
  {\color{soft}\large #3}\par\vspace{5mm}}
\renewcommand{\factsheet}[4]{%
  \begin{tcolorbox}[enhanced, colback=pale, colframe=pale, boxrule=0pt, sharp corners,
    left=7pt, right=7pt, top=5pt, bottom=5pt, fontupper=\small]
  \begin{tabularx}{\linewidth}{@{}>{\leavevmode\bfseries\color{navy}}p{27mm}>{\leavevmode}X@{}}
  In Pollis & #1\\
  Missing in Julia & #2\\
  Package & #3\\
  Tests & #4
  \end{tabularx}
  \end{tcolorbox}}
\renewcommand{\nmcredit}{%
  \par\vfill
  \begin{minipage}[b]{0.78\linewidth}\footnotesize\color{soft}
  Package written and tested with Claude Code (Anthropic) in Pollis.\par
  \textcopyright\ 2026 Antonio Saragga Seabra\enspace·\enspace Licensed under \href{https://creativecommons.org/licenses/by-sa/4.0/}{CC BY-SA 4.0}\enspace·\enspace Draft, \nmdate
  \end{minipage}\hfill
  \begin{minipage}[b]{0.18\linewidth}\raggedleft\pollislogo{9mm}\par\vspace{2pt}{\footnotesize contact@pollis.org}\end{minipage}}


\begin{document}
\emergencystretch=3em

\nmhead{Datasets}
  {PollisDatasets.jl: the datasets of the examples}
  {69 small datasets, from \fn{mtcars} and \fn{iris} to \fn{mroz} and \fn{challenger}, each with the source to
   cite, its licence and every column described, loaded as Tables.jl column tables.}

\factsheet
  {the data of the Code Examples, installed with Pollis; the database \emph{PollisDatasets} in the Databases view,
   the \fn{duckdb>} terminal and the \fn{sql>} mode, one table per dataset, with its descriptions.}
  {a set of datasets whose licences were checked one by one, and recorded, before they came in. RDatasets.jl
   loads the datasets of R's packages, as DataFrames, without the terms of each.}
  {\fn{PollisDatasets.jl} 0.1.0, from GitHub:\newline \fn{Pkg.add(url = "https://github.com/Trumpingtons/PollisDatasets.jl")}.
   It depends on CSV, TOML and Tables only.}
  {70 of 70: every dataset loads with the rows and the columns its entry gives, and an unknown name is an
   error. Six datasets also checked against R's datasets package, to the last digit printed.}

\nmsection{Why a package}
An example is only worth running if its data comes with it. Pollis's Code Examples need datasets that are small,
well known and free to redistribute with an AGPL-3.0 program; a teacher or a student needs to know where each one
comes from and how to cite it. PollisDatasets.jl puts each dataset in \fn{data/<id>.csv}, with its entry in
\fn{datasets.toml}: a title, a domain, a description, the number of rows, the source to cite, where the file was
taken from, the licence the provider states, the outcome of the licence check, and a description of every column,
with its units or codes. Missing values are empty fields in the file and \fn{missing} in Julia.

\nmsub{What it holds}
\begin{valuetable}{@{}>{\raggedright\arraybackslash}p{40mm}>{\raggedright\arraybackslash}X@{}}
\textbf{Domain} & \textbf{Datasets}\\\midrule
Regression & mtcars, anscombe, datasaurus, galton, evals\\
Econometrics & mroz, card\\
Panel econometrics & grunfeld, produc, gasoline, petersen\\
Causal inference & lalonde, cps1, cps2, cps3, mpdta, divorce\\
Count data & horsekicks, insectsprays, biochemists, ships, challenger, assortive\_mating\\
Classification & penguins, iris, titanic, wine, breastwisconsin, glass, sonar, vowel, ionosphere, germancredit,
  image, ecoli, votes, letters, satimage, vehicle, liver, soybean, pima\\
Time series & airpassengers, nile, co2, bg96, dgs10, lynxmackenzie, airquality, sp500, gamestop, gnp82, gnpc96\\
Time series classification & ecg200, earthquakes, twoleadecg\\
Survival & lung, veteran, pbc\\
Bayesian / multilevel & eightschools, radon, sleepstudy\\
Dynamical systems & lynxhare\\
Epidemiology & italycovid\\
Density / mixtures & faithful\\
Market baskets & groceries\\
Text & pollis\_guide, the sections of The Pollis Guide\\
Simulated & sim\_glmm, sim\_clf\\
\end{valuetable}

\nmsection{The licences}
Every dataset passed a licence check before it came in, and its \fn{status} records the outcome:
\begin{valuetable}{@{}>{\raggedright\arraybackslash}p{30mm}>{\raggedright\arraybackslash}X r@{}}
\textbf{Status} & \textbf{Meaning} & \textbf{Datasets}\\\midrule
\fn{open} & an explicit open licence (CC0, CC BY), or permission given & 23\\
\fn{public-domain} & a work of the US government, or published before 1930 & 7\\
\fn{redistributed} & no explicit licence; a small factual table redistributed for decades in CRAN and textbook
  packages & 37\\
\fn{simulated} & made by the Julia code in the entry's \fn{code}, which states the true model; MIT, like the
  package & 2\\
\end{valuetable}
The datasets that failed are listed at the top of \fn{datasets.toml} with the reason: the Fama-French factors (no
grant to redistribute: fetch them live, from \menupath{Explore\gt Access Data Libraries}), the SILSO sunspots (CC
BY-NC, whose non-commercial clause AGPL-3.0 cannot carry) and Boston Housing (an ethically problematic variable).
The S\&P 500 series came in only as a copy dedicated to the public domain (CC0), with the source named.

\nmsection{In Julia}
\begin{valuetable}{@{}>{\raggedright\arraybackslash}p{40mm}>{\raggedright\arraybackslash}X@{}}
\textbf{Function} & \textbf{Returns}\\\midrule
\fn{datasets()} & the catalogue, one row per dataset: \fn{id}, \fn{title}, \fn{domain}, \fn{rows}\\
\fn{dataset("mroz")} & the dataset as a Tables.jl column table, a \fn{NamedTuple} of vectors: \fn{d.wage} is a
  column\\
\fn{datasetinfo("mroz")} & the description, the source to cite, the licence and the columns\\
\end{valuetable}
A column table goes straight into GLM.jl, MLJ.jl, Makie or any Tables.jl reader; Pollis's examples never turn it
into a DataFrame.

\nmsection{In Pollis}
\begin{itemize}
  \item \textbf{Installed with Pollis.} When the Julia REPL starts, Pollis adds PollisDatasets.jl to an environment
    of its own, \fn{\textasciitilde/.pollis/julia/environments}, on the load path, so \fn{using PollisDatasets}
    works in any project without touching its \fn{Project.toml}, and updates it when GitHub has a newer commit.
  \item \textbf{The database.} \fn{PollisDatasets.duckdb}, in \fn{\textasciitilde/.pollis/databases}, has one table
    per dataset, named by its id, with the title, description, source and licence as the table's comment and each
    column's description as its own. It is attached read-only to every connection Pollis makes, so
    \fn{SELECT * FROM PollisDatasets.mtcars} works in the \fn{sql>} mode and the \fn{duckdb>} terminal, and the
    descriptions show in the Databases view. It is rebuilt when the installed package or its files change.
  \item \textbf{New datasets announced.} When a version brings new datasets, the Welcome page names them, with
    \menupath{Open Databases View}.
\end{itemize}

\nmsection{A worked example}
\nmsub{Anscombe's quartet}
\begin{nbox}{}
\ttfamily\small
using PollisDatasets, Statistics\\
datasetinfo("anscombe")\hfill \# 11 rows; Anscombe (1973), The American Statistician\\
a = dataset("anscombe")\\
x, y = a.x1, a.y1\hfill \# and x2, y2 to x4, y4\\
b = cov(x, y) / var(x); (mean(y) - b * mean(x), b, cor(x, y))
\end{nbox}
Four sets of eleven points, built by \nmcite{anscombe}{Anscombe (1973)} to show why a regression should be drawn
before it is believed: the same means, the same variances, the same correlation and the same fitted line, from a
straight line with noise, a curve, a line with one outlier, and a vertical cluster with one far point.
\begin{valuetable}{@{}lrrrrrrr@{}}
\textbf{Set} & \textbf{mean x} & \textbf{var x} & \textbf{mean y} & \textbf{var y} & \textbf{cor} &
  \textbf{intercept} & \textbf{slope}\\\midrule
1 & 9 & 11 & 7.500909 & 4.127269 & 0.816421 & 3.000091 & 0.500091\\
2 & 9 & 11 & 7.500909 & 4.127629 & 0.816237 & 3.000909 & 0.500000\\
3 & 9 & 11 & 7.500000 & 4.122620 & 0.816287 & 3.002455 & 0.499727\\
4 & 9 & 11 & 7.500909 & 4.123249 & 0.816521 & 3.001727 & 0.499909\\
\end{valuetable}
The values are Julia's on \fn{dataset("anscombe")}; R's \fn{lm} on its own \fn{anscombe} gives the same to six
decimals.

\nmsection{What the tests show}
\nmsub{The catalogue and the files agree}
\fn{test/runtests.jl} loads every dataset in the catalogue and checks that it has the number of rows its entry
gives and the columns its entry describes, in the same order; and that \fn{dataset("nope")} is an error that lists
the names there are. 70 of 70 pass, in 7 seconds.

\nmsub{Against R}
Datasets that come from R's datasets package were compared with R 4.5's own copies:
\begin{valuetable}{@{}llrrrr@{}}
\textbf{Dataset} & \textbf{Column} & \textbf{Rows} & \textbf{Values} & \textbf{Sum} & \textbf{Mean}\\\midrule
mtcars & mpg & 32 & 32 & 642.9 & 20.09062\\
iris & sepal\_length & 150 & 150 & 876.5 & 5.843333\\
faithful & eruptions & 272 & 272 & 948.677 & 3.487783\\
airquality & ozone & 153 & 116 & 4887 & 42.12931\\
nile & flow & 100 & 100 & 91935 & 919.35\\
airpassengers & passengers & 144 & 144 & 40363 & 280.2986\\
\end{valuetable}
Both give these values, R to the digits it prints; the 37 ozone readings missing in New York's summer of 1973 are
\fn{missing} in Julia as they are \fn{NA} in R.

\nmsection{Adding a dataset}
In the GitHub clone of PollisDatasets.jl, never in the installed copy: check the licence and choose the
\fn{status}; add \fn{data/<id>.csv} and its entry in \fn{datasets.toml}, with every column described; run the tests;
commit and push. At the next start of the Julia REPL, every Pollis installs the new version, rebuilds the database
and announces the dataset on the Welcome page. \menupath{Help\gt Documentation\gt Drill-down: Data Packages} follows
the whole road.

\nmsection{References}
{\footnotesize\sloppy
\nmref{anscombe}{Anscombe, F. J. (1973). Graphs in statistical analysis. The American Statistician 27(1), 17-21. \nmdoi{10.1080/00031305.1973.10478966}}
\nmref{rcore}{R Core Team. The R datasets package. \url{https://stat.ethz.ch/R-manual/R-devel/library/datasets/html/00Index.html}}
\nmref{software}{Software: CSV.jl, Tables.jl and TOML (MIT). The source of each dataset is in its entry of \fn{datasets.toml}.}
}

\nmcredit
\end{document}
